\documentclass[10pt,a4paper]{amsart}
\usepackage[margin=19mm]{geometry}
\usepackage{amsmath,amssymb,mathtools}
\usepackage[scr=rsfs,cal=euler]{mathalfa}
\usepackage{xfrac}
\usepackage[unicode,hidelinks,bookmarks=false]{hyperref}
\newcommand{\ie}{i.e.\ }
\newcommand{\cf}{cf.\ }
\newcommand{\resp}{respectively\ }
\newcommand{\Subsection}{\subsection}
\providecommand{\nobreakdash}{\nobreak\hbox{-}}
\setlength{\emergencystretch}{2em}
\multlinegap0pt
\usepackage{bm}
\usepackage{bbm}
\usepackage{dsfont}
\usepackage{xspace}
\usepackage{cancel}
\usepackage{mathtools}
\usepackage{amsthm}
\makeatletter
\@ifundefined{lemma}{\newtheorem{lemma}{Lemma}}{}
\@ifundefined{proposition}{\newtheorem{proposition}[lemma]{Proposition}}{}
\makeatother
\newcommand{\bl}[2]{(#1,#2)}
\title[Volumes of consecutively defined sets]{Volumes of consecutively defined sets}
\author{Richard Ehrenborg, Evan Henning}
\date{Source: arXiv:2608.10236v1 (CC BY 4.0)}
\begin{document}
\maketitle

\noindent\emph{Source and licence.} Volumes of consecutively defined sets. Version arXiv:2608.10236v1 (CC BY 4.0), distributed under the Creative Commons Attribution 4.0 International licence, as recorded in the arXiv licence metadata for this exact version and shown on the abstract page https://arxiv.org/abs/2608.10236v1. Full rights notice: licenses/arxiv-2608.10236-CC-BY-4.0.txt.\par\smallskip

\noindent\textit{Editorial scope.} Counted unit: the Proposition whose source label is \texttt{equation\_eigenvalue} in the article (source0.tex lines 604--638, source label \texttt{equation\_eigenvalue}), delivered together with the article's own complete original proof of it (source0.tex lines 639--813, one \texttt{proof} environment of 175 lines). No mathematics is written, repaired or re-derived: statement and proof are the article's own text line for line. Required context retained uncounted: proposition\_general\_2nd\_order (lines 523--541); lemma\_unit\_length (lines 816--827). Every definition, display and label of those lines is retained unchanged. Omitted from this package and not redistributed: 1--522, 542--603, 814--815, 828--1190. Nothing inside a retained excerpt is omitted or altered: every retained body is delivered exactly as the article's lines are written, so no omission receipt of any kind accompanies this package. Citations: the retained excerpts carry no citation command at all, so no bibliography entry of the article is transcribed and no reference is redistributed. Reference adaptation: none was needed and none was made; every reference inside the delivered unit resolves inside it, so no unresolved reference or citation is produced. Printed slips kept literal: no printed word, symbol, hypothesis or number of the counted statement or of its proof was altered. Layout and class: the article's own class is \texttt{amsart} with packages including geometry, amssymb, amsmath, amsthm, color, tikz, enumerate, mathrsfs, pdfpages, hyperref, textcomp, xcolor, listings, fullpage; the wrapper is the lane's pinned \texttt{amsart} editorial wrapper at 10pt on A4, which restates the theorem-like environments the retained text needs and adds the article's own macro definitions that the retained text needs (\texttt{\textbackslash bl}), with extra wrapper packages (bm, bbm, dsfont, xspace, cancel, mathtools). The delivered numbering is the unit's own. Assets: no retained span contains a figure environment, a graphics-inclusion command, a PDF-page inclusion, a table or a TikZ picture; no image file is copied into this package, the package declares no asset dependency and no image file is redistributed with it. Licence: version arXiv:2608.10236v1 of this article is distributed under the Creative Commons Attribution 4.0 International licence, as recorded in the arXiv licence metadata for this exact version and shown on the abstract page https://arxiv.org/abs/2608.10236v1. A full rights notice is delivered at licenses/arxiv-2608.10236-CC-BY-4.0.txt. Characters: every retained excerpt is pure ASCII, so the delivered engine, XeLaTeX with its default Latin Modern text font, reports no missing character.
\begin{proposition}
Assume that $c(t)$ is a piecewise twice differentiable function.
Let $\varphi(t)$ be an eigenfunction
of the operator $T$ with eigenvalue~$\lambda$.
Then the eigenfunction $\varphi(t)$ satisfies
the second order linear differential equation
\begin{align*}
\varphi''(t)
-
\frac{c''(t)}{c'(t)} \cdot \varphi'(t)
-
\frac{c'(t)}{\lambda^{2}} \cdot \varphi(t)
=
0,
\end{align*}
and satisfies the boundary conditions
$\varphi(1) = 0$ and $\varphi'(0) = 0$.
\label{proposition_general_2nd_order}
\end{proposition}

\begin{lemma}
The following two identities hold:
\begin{align*}
\bl{1}{\varphi_{k}}
& = 
\frac{\lambda_{k}}{\sqrt{\alpha}} , &
\bl{\varphi_{k}}{\varphi_{k}}
& = 
1.
\end{align*}
\label{lemma_unit_length}
\end{lemma}

\begin{proposition}
The eigenvalues of the operator $T$ are given by
\begin{align}
\lambda_{k}
& =
\frac{\sqrt{(1-\alpha) \cdot \alpha}}
{
\arctan\left(\sqrt{\frac{1-\alpha}{\alpha}}\right)
+
\pi \cdot k
} ,
\label{equation_eigenvalue}
\end{align}
where $k$ is an integer.
The eigenfunction associated with the eigenvalue $\lambda_{k}$ is
\begin{align*}
\varphi_{k}(t)
& =
\begin{cases}
\frac{1}{\sqrt{\alpha}}
\cdot
\cos\left( \sqrt{\frac{1-\alpha}{\alpha}} \cdot \frac{1}{\lambda_{k}} \cdot  t \right)
&
\text{if } 0 \leq t \leq \alpha , 
\\
\frac{1}{\sqrt{1-\alpha}}
\cdot
\sin\left( \sqrt{\frac{\alpha}{1-\alpha}} \cdot \frac{1}{\lambda_{k}} \cdot (1-t) \right)
&
\text{if } \alpha \leq t \leq 1 . 
\end{cases}
\end{align*}
Furthermore, the algebraic multiplicity of each of the non-zero
eigenvalues is $1$.
\end{proposition}

\begin{proof}
Let $\varphi$ be an eigenfunction,
that is, $\lambda \cdot \varphi = T(\varphi)$.
Assume that $\varphi$ is of the form
\begin{align*}
\varphi(t)
& =
\begin{cases}
p(t) & \text{if } 0 \leq t \leq \alpha , \\
q(t) & \text{if } \alpha \leq t \leq 1 .
\end{cases}
\end{align*}
The derivative of $b(t)$ is given by
\begin{align*}
b'(t)
& = 
\begin{cases}
-(1-\alpha)/\alpha & \text{ for } 0 \leq t \leq \alpha, \\
-\alpha/(1-\alpha) & \text{ for } \alpha \leq t \leq 1 .
\end{cases}
\end{align*}
Note also that $b''(t) = 0$.
Hence Proposition~\ref{proposition_general_2nd_order}
implies that
\begin{align*}
p''(t) 
+
\frac{1}{\lambda^{2}} \cdot \frac{1-\alpha}{\alpha} \cdot p(t) 
& =
0 ,
&&
0 \leq t \leq \alpha,
\\
q''(t)
+
\frac{1}{\lambda^{2}} \cdot \frac{\alpha}{1-\alpha} \cdot q(t) 
& =
0 ,
&&
\alpha \leq t \leq 1.
\end{align*}
These two second order linear differential equations
have the general solutions
\begin{align}
p(t)
& =
A \cdot \cos\left( \sqrt{\frac{1-\alpha}{\alpha}} \cdot \frac{1}{\lambda} \cdot  t \right)
+
B \cdot \sin\left( \sqrt{\frac{1-\alpha}{\alpha}} \cdot \frac{1}{\lambda} \cdot t \right) ,
\label{equation_p_explicit_A_B}
\\
q(t)
& =
C \cdot \cos\left( \sqrt{\frac{\alpha}{1-\alpha}} \cdot \frac{1}{\lambda} \cdot  t \right)
+
D \cdot \sin\left( \sqrt{\frac{\alpha}{1-\alpha}} \cdot \frac{1}{\lambda} \cdot t \right) .
\label{equation_q_explicit_C_D}
\end{align}
It remains to find the four constants $A$, $B$, $C$ and $D$
and an equation for the eigenvalue $\lambda$.

Note that $q(1) = \varphi(1) = 0$
and $p'(0) = \varphi'(0) = 0$
by
Proposition~\ref{proposition_general_2nd_order}.
This last observation implies $B=0$.
Now see that $A = 0$
would imply that $\varphi(t) = 0$.
Hence we obtain that the constant $A$ is non-zero.
Furthermore, setting $t=1$ in
equation~\eqref{equation_q_explicit_C_D}
yields
\begin{align*}
\tan\left( \sqrt{\frac{\alpha}{1-\alpha}} \cdot \frac{1}{\lambda} \right)
& =
- \frac{C}{D} .
\end{align*}
Thus we may select
\begin{align*}
C
& =
\frac{1}{\sqrt{1-\alpha}}
\cdot
\sin\left( \sqrt{\frac{\alpha}{1-\alpha}} \cdot \frac{1}{\lambda} \right),
&
D
& =
- \frac{1}{\sqrt{1-\alpha}}
\cdot
\cos\left( \sqrt{\frac{\alpha}{1-\alpha}} \cdot \frac{1}{\lambda} \right) .
\end{align*}
We are selecting the coefficients such that the eigenfunction will have
unit length; see Lemma~\ref{lemma_unit_length}.
Now $q(t)$ has the expression
\begin{align}
q(t)
& =
\frac{1}{\sqrt{1-\alpha}}
\cdot
\sin\left( \sqrt{\frac{\alpha}{1-\alpha}} \cdot \frac{1}{\lambda} \cdot (1-t) \right) .
\label{equation_q_explicit_again}
\end{align}


Since $\varphi(t)$ is continuous we have
the condition $p(\alpha) = q(\alpha)$.
From equations~\eqref{equation_p_explicit_A_B}
and~\eqref{equation_q_explicit_again}
we have
\begin{align*}
A \cdot \cos\left( \sqrt{(1-\alpha) \cdot \alpha} \cdot \frac{1}{\lambda} \right)
& =
\frac{1}{\sqrt{1-\alpha}}
\cdot
\sin\left( \sqrt{(1-\alpha) \cdot \alpha} \cdot \frac{1}{\lambda} \right) .
\end{align*}
Hence the constant $A$ is given by the tangent function
\begin{align*}
A
& =
\frac{1}{\sqrt{1-\alpha}}
\cdot
\tan\left( \sqrt{(1-\alpha) \cdot \alpha} \cdot \frac{1}{\lambda} \right) .
\end{align*}
Next we compute the integral
\begin{align*}
\int_{0}^{\alpha} p(t) \: dt
& =
A \cdot 
\int_{0}^{\alpha} 
\cos\left( \sqrt{\frac{1-\alpha}{\alpha}} \cdot \frac{1}{\lambda} \cdot  t \right)
\: dt
\\
& =
A \cdot 
\sqrt{\frac{\alpha}{1-\alpha}} \cdot \lambda \cdot
\sin\left( \sqrt{(1-\alpha) \cdot \alpha} \cdot \frac{1}{\lambda} \right) .
\end{align*}
Since
$\lambda \cdot p(\alpha)
= \lambda \cdot \varphi(\alpha)
= \int_{0}^{\alpha} \varphi(t) \: dt
= \int_{0}^{\alpha} p(t) \: dt$
we have
\begin{align*}
\lambda \cdot
A \cdot \cos\left( \sqrt{(1-\alpha) \cdot \alpha} \cdot \frac{1}{\lambda} \right)
& =
A \cdot 
\sqrt{\frac{\alpha}{1-\alpha}} \cdot \lambda \cdot
\sin\left( \sqrt{(1-\alpha) \cdot \alpha} \cdot \frac{1}{\lambda} \right).
\end{align*}
Canceling the factor $A \cdot \lambda$
we obtain
\begin{align*}
\sqrt{\frac{1-\alpha}{\alpha}}
& =
\tan\left( \sqrt{(1-\alpha) \cdot \alpha} \cdot \frac{1}{\lambda} \right) .
\end{align*}
This yields a more explicit expression
$1/\sqrt{\alpha}$ for the constant $A$.
Furthermore we can now solve for the eigenvalue $\lambda$
\begin{align*}
\arctan\left(\sqrt{\frac{1-\alpha}{\alpha}}\right)
+
\pi \cdot k
& =
\sqrt{(1-\alpha) \cdot \alpha} \cdot \frac{1}{\lambda} .
\end{align*}
Solving for $\lambda$ yields~\eqref{equation_eigenvalue}.
Finally, since we obtain a unique eigenfunction,
the geometric multiplicity is $1$ of the non-zero eigenvalues.
Since there are no generalized eigenfunctions,
the algebraic multiplicity is the same as the geometric multiplicity.
\end{proof}

\end{document}
